\chapter[\hspace{0.6em}标准化与参数 g 公式]{标准化与参数 g 公式}

本章说明如何用标准化估计戒烟对体重增加的平均因果效应。所用的观测数据集与上一章相同。标准化在第 2 章已经介绍过，但当时只把它当作一种非参数方法。现在我们说明如何把模型与标准化结合起来使用，从而处理协变量众多、治疗非二值的高维问题。与上一章一样，我们提供执行这些分析的计算机代码。

在实践中，研究者常常要在 IP 加权与标准化之间选择一种，作为从观测数据获得效应估计的分析途径。两种方法依据相同的可识别条件，但依赖的建模假设不同。

\section{标准化作为 IP 加权的替代}

上一章我们用 IP 加权估计了戒烟 $A$（1：是，0：否）对体重增加 $Y$（以 kg 计）的平均因果效应。本章我们用标准化估计同一效应。分析同样基于 NHEFS 数据：1629 名 25 至 74 岁的吸烟者，均接受过一次基线访视与约 10 年后的随访。其中 1566 人在随访时测量了体重，因而未被删失（$C=0$）。

我们定义 $\E[Y^{a,c=0}]$ 为：若所有人都接受治疗水平 $a$、且无人被删失，本可观测到的平均体重增加。戒烟的平均因果效应可表示为差值 $\E[Y^{a=1,c=0}]-\E[Y^{a=0,c=0}]$，即“人人受治疗且未删失”与“人人未受治疗且未删失”两种情形下本可观测到的平均体重之差。

如表 12.1 所示，戒烟者（$A=1$）与未戒烟者（$A=0$）在体重增加的各预测变量的分布上存在差异。观测到的关联差 $\E[Y\given A=1,C=0]-\E[Y\given A=0,C=0]=2.5$，预期不同于因果差 $\E[Y^{a=1,c=0}]-\E[Y^{a=0,c=0}]$。我们同样假定变量向量 $L$ 足以校正混杂与选择偏倚，且 $L$ 包含以下基线变量：性别（0：男，1：女）、年龄（岁）、种族（0：白人，1：其他）、教育（5 类）、吸烟的强度与年限（每日吸烟支数与吸烟年数）、日常生活中的体力活动（3 类）、休闲锻炼（3 类）以及体重（kg）。

校正变量 $L$ 的一种方法是 IP 加权：它构造一个伪总体，使 $L$ 中各变量的分布在受治疗者与未受治疗者中相同。于是，在给定 $L$ 的可交换性与正值性假设下，我们把伪总体中的平均结局 $\widehat{\E}[Y\given A=a,C=0]$ 直接算出来，作为 $\E[Y^{a,c=0}]$ 的估计。若 $A$ 是连续治疗（与我们的例子不同），我们还需要一个结构模型，以便对 $A$ 的一切可能取值估计伪总体中的 $\E[Y\given A,C=0]$。\mnote{与上一章一样，我们假定为校正 $C$ 所需的那部分 $L$ 不受 $A$ 影响。否则，我们就需要用第三部分所述的更一般的方法。}IP 加权需要估计治疗与删失的联合分布。对于二值的戒烟治疗，我们估计了 $\Pr[A=a,C=0\given L]$，并用这一联合概率作分母计算 IP 概率权重。

正如第 2 章所讨论的，IP 加权的替代方法是标准化。在给定 $L$ 中变量的可交换性与正值性下，未删失受治疗者的标准化平均结局，是“人人受治疗且保持未删失”时平均结局 $\E[Y^{a=1,c=0}]$ 的一致估计量。类似地，未删失未受治疗者的标准化平均结局，是“人人未受治疗且保持未删失”时平均结局 $\E[Y^{a=0,c=0}]$ 的一致估计量。关于 IP 加权与标准化中偏离正值性影响的相对大小，见精要 13.1。

\renewcommand{\FPname}{精要（Fine Point）}
\begin{finepoint}[结构性正值性]
缺乏结构性正值性，会使 IP 加权无法识别整个总体的平均因果效应。标准化同样需要正值性，因为当 $\Pr[A=a\given L=l]=0$ 而 $\Pr[L=l]\neq 0$ 时，条件平均结局 $\E[Y\given A=a,L=l]$ 无定义。

但偏离正值性对 IP 加权估计与依赖参数模型的标准化估计的实际影响，可能差别很大。使用标准化时，只要愿意依赖参数外推，就可以忽略正值性的缺失。也就是说，可以为 $\E[Y\given A,L]$ 拟合一个模型，让它把结构性零的那些层平滑掉。这种平滑会给估计引入偏倚，因此估计值周围名义上的 95\% 置信区间覆盖真实效应的频率将低于 95\%。还要注意，在结构性正值性成问题的情形下，建模的目的有所不同：我们建模，不是因为数据不够，而是因为要估计一个即便有无限数据也无法识别的量（由于结构性非正值性）。这是一个重要的区别。

一般而言，在正值性被违背或接近被违背时，标准化的治疗效应标准误小于 IP 加权。这并不必然意味着标准化优于 IP 加权；偏倚上的差异可能淹没标准误上的差异。
\end{finepoint}
\renewcommand{\FPname}{精要}

要计算未删失受治疗者的标准化平均结局，首先需要计算各混杂因子层 $l$ 中未删失受治疗者的平均结局，即在每一层 $l$ 中的条件均值 $\E[Y\given A=1,C=0,L=l]$。在我们的戒烟例子中，需要计算戒烟且保持未删失者中，在由 $L$ 中 9 个变量取值组合所定义的（可能多达数百万个）各层内，平均体重增加 $Y$ 是多少。\mnote{技术注记 2.3 证明：在条件可交换性、正值性与一致性下，受治疗者的标准化平均等于人人受治疗时的均值。推广到删失情形是平凡的：只需在证明与定义中把 $a$ 换成 $(A=a,C=0)$。}

未删失受治疗者的标准化均值，是这些条件均值的加权平均，权重为各取值 $l$ 在研究总体中的占比，即 $\Pr[L=l]$。也就是说，个体数最多的那一层的条件均值，在标准化均值的计算中权重最大。未删失未受治疗者的标准化均值算法类似，只需把条件事件中的 $A=1$ 换成 $A=0$。

更紧凑地，接受治疗水平 $a$ 的未删失者的标准化均值为
\[
\sum_l \E[Y\given A=a,C=0,L=l]\times\Pr[L=l]
\]
当 $L$ 中某些变量如本例那样是连续的，就需要把 $\Pr[L=l]$ 换为概率密度函数（pdf）$f_L(l)$，上述求和变为积分。\mnote{受治疗者的平均因果效应可按技术注记 4.1 所述用标准化估计。只需把上式中的 $\Pr[L=l]$ 换成 $\Pr[L=l\given A=1]$。}

接下来两节说明如何估计结局 $Y$ 的条件均值与混杂因子 $L$ 的分布——这是估计标准化均值所需的两类量。

\section{用模型估计平均结局}

理想情况下，我们会非参数地估计这组条件均值 $\E[Y\given A=1,C=0,L=l]$。我们会在由 $L$ 中各变量不同取值组合所定义的每一层内，计算未删失受治疗者的平均结局。这正是 2.3 节所做的：那里计算所需的全部信息都取自表 2.2。

但是，当数据像本例这样高维、混杂因子众多且其中一些有多个水平时，非参数估计 $\E[Y\given A=1,C=0,L=l]$ 根本不可行。当只有 403 名受治疗者、却分散在数百万个层中时，我们无法得到有意义的、分层的受治疗者平均结局非参数估计。我们只能求助于建模。同样的道理适用于未删失未受治疗者的条件平均结局 $\E[Y\given A=0,C=0,L=l]$。

为在由 $L$ 定义的数百万个层中各得到一个 $\E[Y\given A=a,C=0,L=l]$ 的参数估计，我们拟合一个平均体重增加的线性回归模型，以治疗 $A$ 与 $L$ 中全部 9 个混杂因子为协变量。对（准）连续协变量年龄、体重、吸烟强度与吸烟年限，我们使用了线性项与二次项。也就是说，我们的模型对 $\E[Y\given A=a,C=0,L=l]$ 的可能取值施加了限制：连续协变量与平均结局之间的条件关系可以用一条抛物线表示。我们还加入了戒烟 $A$ 与吸烟强度之间的乘积项。也就是说，我们的模型施加的限制是：各协变量对均值的贡献不依赖于其他协变量的贡献，唯一的例外是戒烟 $A$ 的贡献随既往吸烟强度线性变化。\mnote{代码：程序 13.1}

在这些参数限制下，我们得到了 $A$ 与 $L$ 每一取值组合的估计 $\widehat{\E}[Y\given A=a,C=0,L=l]$，因而也得到研究总体中 403 名未删失受治疗者（$A=1,C=0$）与 1163 名未删失未受治疗者（$A=0,C=0$）各自的估计值。

例如，我们估计取值组合为\{未戒烟者、男性、白人、26 岁、大学辍学、每日 15 支烟、吸烟史 12 年、中等锻炼、体力活动很活跃、体重 112 kg\}的个体，其平均体重增加为 0.34 kg（唯一标识为 24770 的个体恰好具有这组取值，你可以看看他的预测值）。
\[
\int \E[Y\given A=a,C=0,L=l]\,dF_L(l)
\]
总体而言，估计的体重增加均值为 2.6 kg，与观测到的体重增加均值相同；后者在不同协变量组合之间从 $-41.3$ kg 到 48.5 kg 不等。

请记住，我们的目标是估计受治疗者（$A=1$）与未受治疗者（$A=0$）中的标准化均值 $\sum_l \E[Y\given A=a,C=0,L=l]\times\Pr[L=l]$。更形式化地说，标准化均值应写成积分，因为 $L$ 中有些变量本质上是连续的，其分布无法用概率函数表示。撇开这些记号问题，我们已经对治疗 $A$ 与混杂因子 $L$ 的一切取值估计出了均值 $\E[Y\given A=a,C=0,L=l]$。\mnote{一般而言，$Y$ 的标准化均值写作 $\int \E[Y\given A=a,C=0,L=l]\,dF_L(l)$，其中 $F_L(\cdot)$ 是 $L$ 中随机变量的联合累积分布函数（cdf）。当本章这样、$L$ 是不受治疗影响的基线协变量向量时，我们可以对 $L$ 的观测值取平均，从而非参数地估计这一积分。}

\section{把平均结局标准化到混杂因子分布}

标准化均值是条件均值 $\E[Y\given A=a,C=0,L=l]$ 的加权平均。当 $L$ 中所有变量都是离散的，每个均值所得的权重等于取值 $L=l$ 的个体所占比例，即 $\Pr[L=l]$。原则上，这些比例 $\Pr[L=l]$ 可以从数据中非参数地算出来：把由 $L=l$ 定义的层中的个体数除以总体中的个体总数。这正是 2.3 节所做的：那里计算所需的全部信息都取自表 2.2。然而，对于混杂因子众多、其中一些有多个水平的高维数据，如我们的戒烟例子，这种方法变得繁琐。

幸运的是，我们并不需要估计 $\Pr[L=l]$。我们只需对研究中每个个体 $i$ 的 $l$ 值估计 $\E[Y\given A=a,C=0,L=l]$，然后计算平均
\[
\frac{1}{n}\sum_{i=1}^{n}\widehat{\E}[Y\given A=a,C=0,L_i],
\]
其中 $n$ 是研究中的个体数。之所以可以这样，是因为加权平均 $\sum_l \E[Y\given A=a,C=0,L=l]\Pr[L=l]$ 也可写成二重期望 $\E\bigl[\E[Y\given A=a,C=0,L]\bigr]$。

下面介绍一种简单的计算方法，用于在混杂因子众多时估计受治疗者（$A=1$）与未受治疗者（$A=0$）的标准化均值 $\sum_l \E[Y\given A=a,C=0,L=l]\times\Pr[L=l]$，而始终不必显式估计 $\Pr[L=l]$。我们先把该方法用于表 2.2 的数据——那里没有删失，混杂因子 $L$ 只有一个二水平变量，$Y$ 是二值结局，即均值 $\E[Y\given A=a,C=0,L=l]$ 就是发生结局的风险 $\Pr[Y=1\given A=a,L=l]$。然后再把它用于有删失、混杂因子众多的真实数据。该方法有 4 步：扩展数据集、结局建模、预测、取平均完成标准化。

表 2.2 有 20 行，每个研究个体一行。现在我们新建一个数据集，把表 2.2 的数据复制三份。也就是说，分析数据集有 60 行，分成三个各含 20 人的区块。第一区块的 20 行原样保留，即与表 2.2 的数据完全相同。我们按下表所示修改第二、三区块的数据。在第二区块中，把全部 20 人的 $A$ 值设为 0（未治疗）；在第三区块中，把所有人的 $A$ 值设为 1（治疗）。在第二、三区块中，我们删除所有个体的结局数据，即把变量 $Y$ 赋为缺失值。如下文所述，我们将用第二区块估计未受治疗者的标准化均值，用第三区块估计受治疗者的标准化均值。

\begin{center}
\small
\begin{tabular}{lccc@{\hspace{2.6em}}lccc}
\toprule
\multicolumn{4}{c}{第二区块：全部未治疗} & \multicolumn{4}{c}{第三区块：全部治疗}\\
\cmidrule(lr){1-4}\cmidrule(lr){5-8}
 & $L$ & $A$ & $Y$ & & $L$ & $A$ & $Y$\\
\midrule
Rheia      & 0 & 0 & . & Rheia      & 0 & 1 & .\\
Kronos     & 0 & 0 & . & Kronos     & 0 & 1 & .\\
Demeter    & 0 & 0 & . & Demeter    & 0 & 1 & .\\
Hades      & 0 & 0 & . & Hades      & 0 & 1 & .\\
Hestia     & 0 & 0 & . & Hestia     & 0 & 1 & .\\
Poseidon   & 0 & 0 & . & Poseidon   & 0 & 1 & .\\
Hera       & 0 & 0 & . & Hera       & 0 & 1 & .\\
Zeus       & 0 & 0 & . & Zeus       & 0 & 1 & .\\
Artemis    & 1 & 0 & . & Artemis    & 1 & 1 & .\\
Apollo     & 1 & 0 & . & Apollo     & 1 & 1 & .\\
Leto       & 1 & 0 & . & Leto       & 1 & 1 & .\\
Ares       & 1 & 0 & . & Ares       & 1 & 1 & .\\
Athena     & 1 & 0 & . & Athena     & 1 & 1 & .\\
Hephaestus & 1 & 0 & . & Hephaestus & 1 & 1 & .\\
Aphrodite  & 1 & 0 & . & Aphrodite  & 1 & 1 & .\\
Polyphemus & 1 & 0 & . & Polyphemus & 1 & 1 & .\\
Persephone & 1 & 0 & . & Persephone & 1 & 1 & .\\
Hermes     & 1 & 0 & . & Hermes     & 1 & 1 & .\\
Hebe       & 1 & 0 & . & Hebe       & 1 & 1 & .\\
Dionysus   & 1 & 0 & . & Dionysus   & 1 & 1 & .\\
\bottomrule
\end{tabular}
\end{center}

接下来，我们用这个三区块数据集拟合一个以治疗 $A$ 与混杂因子 $L$ 为自变量的平均结局回归模型。我们加入乘积项 $A\times L$ 使模型饱和。注意，只有数据集第一区块的行（真实数据）会对模型参数的估计有贡献，因为第二、三区块所有行的结局都是缺失的。

下一步是用第一区块得到的参数估计，预测第二、三区块所有行的结局值。（即把 $L$ 与 $A$ 两列的取值与回归估计结合起来，插补结局 $Y$ 的缺失值。）第二区块的预测结局值，是各 $L$ 取值组合在 $A=0$ 下的均值估计；第三区块的预测值，是各 $L$ 取值组合在 $A=1$ 下的均值估计。

最后，我们计算第二区块中所有预测值的平均。由于 60\% 的行取 $L=1$、40\% 的行取 $L=0$，这个平均值给 $L=1$ 的行更大权重。也就是说，第二区块中所有预测值的平均，恰好就是未受治疗者的标准化平均结局。到此完成。要估计受治疗者的标准化平均结局，我们计算第三区块中所有预测值的平均。\mnote{代码：程序 13.2}

上述步骤得到的标准化均值估计（两者都是 0.5），与 2.3 节的直接计算完全相同。两种做法都完全是非参数的。本章我们没有直接估计 $L$ 的分布，而是对 $L$ 的观测值——即其经验分布——取平均。

在更现实的例子中，如我们的戒烟研究，$L$ 是高维的，用经验分布做标准化才是可行之道。我们研究的步骤与上述表 2.2 数据的步骤类似。我们把第二、三区块加入数据集，按上一节所述拟合 $\E[Y\given A=a,C=0,L=l]$ 的回归模型，并生成预测值。第二区块的平均预测值——未受治疗者的标准化均值——为 1.66；第三区块的平均预测值——受治疗者的标准化均值——为 5.18。\mnote{代码：程序 13.3}因此，我们对因果效应 $\E[Y^{a=1,c=0}]-\E[Y^{a=0,c=0}]$ 的估计为 $5.18-1.66=3.5$ kg。为得到这一估计的 95\% 置信区间，我们使用了称为自助法（bootstrapping）的统计技术（见技术注记 13.1）。\mnote{代码：程序 13.4}总之，我们估计戒烟使体重平均增加 3.5 kg（95\% 置信区间：2.6, 4.5）。

\renewcommand{\TPname}{技术注记（Technical Point）}
\begin{techpoint}[自助法]
在第 10 章，我们讨论了因果效应随机变异的基础。效应估计需要与变异度量一起呈现，例如标准误（或标准误的函数，如 95\% 置信区间）。由于精确估计在计算上有困难，实践中标准误估计常基于大样本近似，而大样本近似依赖渐近论证。然而，有时连大样本近似都过于复杂而无法计算。

自助法是估计标准误与计算 95\% 置信区间的另一种方法。下面勾勒其最简单的版本——非参数自助法，我们用它计算了戒烟效应估计的 95\% 置信区间。

取 1629 人的研究总体。从研究总体中有放回地抽取 1629 人，这样某些原有个人可能出现多次，另一些可能完全不被抽中。这个容量为 1629 的新样本称为一个“自助样本”。在自助样本中计算所关心的效应（例如用正文所述的标准化）。然后再抽取 1629 人、构造第二个自助样本。用与第一个自助样本相同的方法在第二个自助样本中计算所关心的效应。由于随机性，第一、第二个自助样本一般会包含每个个体的不同份数，因此得到不同的效应估计。把这一过程重复大量次数（比如 1000 个自助样本）。结果是，1000 个自助样本效应估计的标准差，一致地估计研究总体中效应估计的标准误。然后用通常的正态近似计算 95\% 置信区间：估计的标准误乘以 $\pm 1.96$。关于自助法背后的统计理论导引，参见 Wasserman (2004)。

我们用这一自助法、取 1000 个自助样本，得到正文所述标准化均值差的 95\% 置信区间。自助法是一种适用于大样本的一般方法：一般地，当估计量的极限分布为正态时，以该估计量为中心、用非参数自助法估计标准误的 95\% Wald 置信区间，在大样本下是校准的。因此，若标准误用自助法估计，边缘结构模型的 IP 加权估计的 95\% Wald 置信区间将是校准的；但若用前面所述的（稳健方差估计量的平方根）来估计，则常常偏保守、偏宽。

虽然非参数自助法简单，但对很大的数据集可能计算量很大。因此，常见已发表的估计只基于 200 至 500 个自助样本。虽然在我们这个例子里减少样本数会得到几乎相同的置信区间，但情况未必总是如此。要在保留自助法优点的同时更好地应对这些计算挑战，一个巧妙的方法是所谓的“小自助袋”（bag of little bootstraps，Kleiner et al. 2014）。
\end{techpoint}
\renewcommand{\TPname}{技术注记}

\section{IP 加权还是标准化？}

至此我们已经介绍了用模型估计治疗平均因果效应的两种方式：IP 加权（上一章）与标准化（本章）。在戒烟例子中，两者给出的效应估计几乎完全相同。事实上，技术注记 2.3 已证明标准化均值等于 IP 加权均值。

既然上一章已经估计了 IP 加权均值，本章为何还要费力估计标准化均值？答案是：只有在不使用模型估计它们时，IP 加权均值与标准化均值才严格相等。否则，二者预期会有差异。要看清这一点，考察实施 IP 加权或标准化所需建模的量。IP 加权对 $\Pr[A=a,C=0\given L]$ 建模，上一章我们通过为 $\Pr[A=a\given L]$ 与 $\Pr[C=0\given A=a,L]$ 拟合参数 logistic 回归模型来估计它。标准化对条件均值 $\E[Y\given A=a,C=0,L=l]$ 建模，本章我们用参数线性回归模型来估计它。

在实践中，任何模型都不可避免地存在某种程度的误设，而模型误设会引入一定偏倚。但治疗模型（IP 加权）与结局模型（标准化）的误设，一般不会在效应估计中造成方向和大小相同的偏倚。因此，由于不可避免的模型误设会以不同方式影响点估计，IP 加权估计一般不同于标准化估计。两者差别很大，提示我们至少其中一个估计存在严重的模型误设。差别很小并不能保证没有严重误设，但会让人安心——两个模型的误设程度都很严重、且造成的偏倚大小与方向恰好相似，虽在逻辑上可能，却不太会发生。

在戒烟例子中，IP 加权估计与标准化估计相近。四舍五入到一位小数，无论我们拟合治疗 $A$ 的模型（IP 加权）还是结局 $Y$ 的模型（标准化），戒烟所致的体重增加估计都是 3.5 kg。两种情形下我们都没有对混杂因子 $L$ 建模，因为获得 IP 加权估计不需要 $L$ 的分布，而计算标准化估计时我们可以使用 $L$ 的经验分布（一种非参数方法）。

IP 加权与标准化都是 g 公式的估计量；g 公式是 1986 年首次提出的因果推断的一般方法。\mnote{Robins (1986) 描述了把标准化推广到随时间变化的治疗与混杂因子，并将其命名为\emph{g 计算算法公式}。由于这个名字很长，有些作者把它缩写为 g-formula，另一些作者缩写为 g-computation。尽管 g-formula 与 g-computation 是同义词，本书只用前一个说法，因为后一个说法常与第 14 章所述的另一种方法——g 估计——相混淆。}（第三部分给出治疗随时间变化情形下 g 公式的定义。）我们说标准化是一种插入式（plug-in）g 公式估计量，因为它只是把 g 公式中的条件平均结局换成其估计值。当这些估计值如本章这样来自参数模型时，我们把该方法称为\emph{参数 g 公式}。由于这里我们只关心平均因果效应，我们参数地估计了条件平均结局。

更一般地，概率密度函数（pdf）的参数 g 公式需要估计结局在 $A$ 与 $L$ 各水平内的条件分布，以计算其标准化值。在没有随时间变化的混杂因子时（见第三部分），参数 g 公式不需要对混杂因子的分布做参数建模。

通常不必在 IP 加权与参数 g 公式之间做选择。当两种方法都能用来估计因果效应时，就两种都用。此外，只要可能，就使用把治疗模型与结局模型结合在同一估计量中的双重稳健方法。在给定 $L$ 的可交换性与正值性下，只要治疗模型或结局模型之一正确（且无需知道哪一个正确），双重稳健估计量就一致地估计平均因果效应。精要 13.2 讨论了一种特定的双重稳健估计量——双重稳健插入式估计量。技术注记 13.2 讨论了第二种双重稳健估计量——增广 IP 加权估计量。技术注记 13.3 讨论两者之间的数学关系。

最后，请注意我们用参数 g 公式估计的是整个目标总体的平均因果效应。若我们关心的是总体某一特定子集的平均因果效应，可以把计算限制在该子集内。例如，若我们关心性别造成的潜在效应修饰，就会分别估计男性与女性的标准化均值。IP 加权与参数 g 公式都可用于估计整个总体或某一子集的平均因果效应。

\section{我们该多认真地对待这些估计？}

本书第一部分回顾了平均因果效应的定义、估计它所需的假设，以及许多潜在的偏倚。那些讨论纯属概念性的，数据例子也过于简单。一个关键信息是：对观测数据做因果分析时，若能明确地模拟一个（假想的）随机化实验——目标试验——分析会更锐利。

本章与上一章的分析，是我们从真实数据估计因果效应的第一次尝试。使用 IP 加权与参数 g 公式，我们估计：若人人都戒烟，平均体重增加本会是 5.2 kg；若无人戒烟，则为 1.7 kg。两种方法都估计，在这个特定总体中，戒烟平均使体重增加 3.5 kg（95\% 置信区间：2.5, 4.5）。在接下来的几章我们会看到，使用 g 估计、结局回归与倾向性评分会得到相似的估计。

各方法的估计彼此相容，这令人安心，因为每种方法的估计都基于不同的建模假设。然而，观测数据的效应估计总是难逃严厉批评。即便我们不想把效应估计迁移到其他总体（第 4 章），即便个体之间没有干扰，要让我们对目标总体的估计有效，仍需满足许多条件。我们把这些条件归为三组。

第一，可交换性、正值性与一致性这些可识别条件（第 3 章）必须成立，观测研究才像目标试验。戒烟者与未戒烟者必须在给定 9 个已测协变量 $L$ 的条件下可交换（见精要 14.2）。未测混杂（第 7 章）或选择偏倚（第 8 章、精要 12.2）会破坏条件可交换性。正值性要求戒烟者中协变量 $L$ 的分布与未戒烟者中的分布完全重叠。精要 13.1 讨论了偏离正值性对非参数 IP 加权与标准化的不同影响。关于一致性，注意戒烟（如逐渐戒、突然戒）与不戒烟（如每日吸烟量增加 2 支、减少但未减到零）都有多个版本。我们的效应估计对应于目标总体中一个颇为含糊的假想干预：它按这些治疗版本在研究总体中实际出现的频率随机分配它们。其他假想干预可能得到不同的效应估计。\mnote{当估计值被用于指导公共政策或临床医学的决策时，数值估计对确切干预的这种依赖很重要（Hernán 2016）。}

第二，分析中使用的所有变量都必须正确测量。治疗 $A$、结局 $Y$ 或混杂因子 $L$ 的测量误差一般会导致偏倚（第 9 章）。实践中，大多数变量的某种程度的误测不可避免。

第三，分析中使用的所有模型都必须正确设定（第 11 章）。设想治疗模型中连续协变量年龄的正确函数形式不是我们所用的抛物线，而是某个复杂多项式所表示的曲线。那么，即便所有混杂因子都被正确测量并纳入 $L$，IP 加权也不能完全校正混杂。模型误设的影响与混杂因子的测量误差类似。

确保上述每一条件至少近似成立，是研究者最重要的任务。若能保证这些条件成立，数据分析就变得平凡。问题当然在于，永远不能指望这些条件完美成立。未测混杂因子、不重叠的混杂因子分布、定义不良的干预、误测的变量、误设的模型，通常会潜伏在我们的估计背后。其中一些问题可以经验地处理，另一些则仍需依赖实质判断，因而会招致数据无法反驳的批评。例如，我们可以提出不同的模型设定，但无法校正未被测量的变量。\mnote{我们的因果推断要成立，需要以下条件：
\begin{itemize}[leftmargin=1.1em,itemsep=0.05em,topsep=0.15em]
\item 可交换性
\item 正值性
\item 一致性
\item 无测量误差
\item 无模型误设
\end{itemize}}

上述报告的效应估计，只有在这些（苛刻的）条件全部成立时，才是戒烟平均因果效应的无偏估计。我们的研究偏离这些条件越远，效应估计的偏倚可能越大。这些条件无法经验地检验，因为我们缺少反事实结局分布的数据。因此，实践中我们假定上述条件近似满足。这一假定需要专家知识的支持，正如 7.6 节针对混杂所致可交换性缺失所讨论的那样。

然而，专家知识是不完备的。因此，现有专家知识通常与一系列条件相容：从因所有已知混杂因子都已被测量而近乎完全可交换，到因我们或许不知道某些混杂因子而中度缺乏可交换性。所以，实践中我们需要在不同假设下做分析，以考察效应估计对原假设的敏感性。本书中，我们所说的敏感性分析包括：针对混杂的（见精要 7.1 所引文献），可用阴性结局对照（技术注记 7.5）与 g 估计（精要 14.2）来做；针对选择偏倚的（精要 12.1）；以及针对模型误设的（11.5 节）。效应估计对不可验证条件之依赖的敏感性，也可通过定量偏倚分析（精要 10.2）来考察，有时则借助另一套不可验证条件，例如工具变量估计所需的条件（见第 16 章）。理想情况下，敏感性分析应纳入所有因果推断研究项目。\mnote{存在未测混杂因子时，另一套可识别条件是近端因果推断（技术注记 7.3）与前门准则（技术注记 7.4）。}

对从观测数据得出的因果推断，保持健康的怀疑是必要的。要使这种怀疑富有成效，它必须建立在对我们假设之有效性的专家知识之上。少一些随意的因果推断，关键一步是认识到：讨论应主要围绕上述每一条假设展开。我们对待效应估计的认真程度，取决于我们对待赋予其因果解释所需条件的认真程度。

\renewcommand{\FPname}{精要（Fine Point）}
\begin{finepoint}[双重稳健插入式估计量]
上一章介绍 IP 加权，它要求治疗 $A$ 在给定混杂因子 $L$ 下的模型正确。本章介绍标准化，它要求结局 $Y$ 在给定治疗 $A$ 与混杂因子 $L$ 下的模型正确。那么，有没有一种方法，只要求治疗 $A$ 或结局 $Y$ 之一的模型正确？双重稳健估计正是如此。在通常的可识别假设下，只要两个模型中至少有一个正确（且无需知道哪一个正确），双重稳健估计量就一致地估计因果效应。也就是说，双重稳健估计量给了我们两次做对的机会。

双重稳健估计量有很多类型。这里我们介绍一种（Bang and Robins, 2005），用于二值治疗 $A$ 对结局 $Y$ 的平均因果效应。为简单起见，我们考虑无删失的情形。

要得到平均因果效应的双重稳健估计，首先按上一章所述估计 IP 权重 $W^A=1/f(A\given L)$。然后拟合一个本章所述那样的结局回归模型——带标准联系函数的广义线性模型——对 $\E[Y\given A,L,R]$ 建模，并加入协变量 $R$，其中若 $A=1$ 则 $R=W^A$，若 $A=0$ 则 $R=-W^A$。最后，把每个个体的 $A$ 设为 1、由结局模型取预测值，得到 $A=1$ 下标准化平均结局的估计；再把 $A$ 设为 0 重复一遍，得到 $A=0$ 下标准化平均结局的估计。两估计量之差，即平均因果效应的双重稳健插入式估计量。
\end{finepoint}
\renewcommand{\FPname}{精要}


\renewcommand{\TPname}{技术注记（Technical Point）}
\begin{techpoint}[增广 IP 加权估计量]
设我们有二值治疗 $A$、结局 $Y$，以及满足正值性与可交换性的已测变量向量 $L$（假定一致性成立）。为简单起见，我们考虑估计治疗 $a=1$ 下的反事实平均结局 $\E[Y^{a=1}]$，而非因果效应。那么 $\E[Y^{a=1}]$ 可写成 $\E[b(L)]$，其中 $b(L)=\E[Y\given A=1,L]$；也可写成 $\E\bigl[AY/\pi(L)\bigr]$，其中 $\pi(L)=\Pr[A=1\given L]$。本章我们介绍了插入式 g 公式估计量 $\frac{1}{n}\sum_{i=1}^{n}\hat b(L_i)$，它把条件平均结局换成来自 $b(L)$ 的（譬如线性）参数回归模型的估计，并对研究中全部 $n$ 个个体取平均。上一章我们介绍了 Horvitz--Thompson IP 加权估计量 $\frac{1}{n}\sum_{i=1}^{n}A_iY_i/\hat\pi(L_i)$，它把治疗概率换成来自 $\pi(L)$ 的（譬如 logistic）参数回归模型的估计，并对 $n$ 个个体取平均。若 $\hat b(L)$ 离 $b(L)$ 很远，插入式 g 公式估计量的偏倚会很大；若 $\hat\pi(L)$ 离 $\pi(L)$ 很远，IP 加权估计量的偏倚会很大。

$\E[Y^{a=1}]$ 的双重稳健估计量把结局模型的估计 $\hat b(L)$ 与治疗模型的估计 $\hat\pi(L)$ 适当地结合起来。双重稳健估计量有多种形式，如精要 13.2 中针对平均因果效应所述的那一种。所有双重稳健估计量都涉及用含有治疗模型的函数对结局回归模型作校正，也可以看作是用含有结局回归模型的函数对 Horvitz--Thompson 估计量作校正。例如，考虑 $\E[Y^{a=1}]$ 的如下双重稳健估计量：
\[
\widehat{\E}[Y^{a=1}]_{DR}=\frac{1}{n}\sum_{i=1}^{n}\left[\hat b(L_i)+\frac{A_i}{\hat\pi(L_i)}\left(Y_i-\hat b(L_i)\right)\right],
\]
它也可写成 $\frac{1}{n}\sum_{i=1}^{n}\left[\frac{A_iY_i}{\hat\pi(L_i)}-\left(\frac{A_i}{\hat\pi(L_i)}-1\right)\hat b(L_i)\right]$。正是后一公式促使我们把 $\widehat{\E}[Y^{a=1}]_{DR}$ 称为\emph{增广 IP 加权估计量}。

在可交换性与正值性下，若 $\hat b(L)$ 接近 $b(L)$ 或 $\hat\pi(L)$ 接近 $\pi(L)$，该 $\E[Y^{a=1}]$ 的双重稳健估计量的偏倚就小。具体而言，差值 $\widehat{\E}[Y^{a=1}]_{DR}-\E[Y^{a=1}]$ 将依概率收敛到
\[
\E\left[\pi(L)\left(\frac{1}{\hat\pi(L)}-\frac{1}{\pi^*(L)}\right)\left(b(L)-b^*(L)\right)\right],
\]
其中 $\pi^*(l)$ 与 $b^*(l)$ 分别是 $\hat\pi(l)$ 与 $\hat b(l)$ 的概率极限。由此可知，当参数结局模型正确［即 $b^*(l)=b(l)$］或参数治疗模型正确［即 $\pi^*(l)=\pi(l)$］时，我们的双重稳健估计量是（渐近）无偏的。此外，我们无需知道两个模型中哪一个正确。当然，若向量 $L$ 维数很高，就不能指望任何参数模型被正确设定，因而即便我们的双重稳健估计量的偏倚也可能很大。

然而，所有双重稳健估计量都有这样一个性质：偏倚取决于 $\frac{1}{\pi(l)}$ 的估计误差 $\frac{1}{\hat\pi(l)}-\frac{1}{\pi(l)}$ 与 $b(l)$ 的估计误差 $b(l)-\hat b(l)$ 的乘积。正如我们在第 18 章所讨论的，这一性质——称为\emph{二阶偏倚}——使我们能够用机器学习估计量而非标准参数模型来估计 $\pi(l)$ 与 $b(l)$，从而构造偏倚可能很小的 $\E[Y^{a=1}]$ 双重稳健估计量。这是因为在数据量很大的高维情形下，基于复杂算法的机器学习估计量能比标准参数模型给出更准确的 $\pi(l)$ 与 $b(l)$ 估计。
\end{techpoint}
\renewcommand{\TPname}{技术注记}

\renewcommand{\TPname}{技术注记（Technical Point）}
\begin{techpoint}[增广 IP 加权估计量与双重稳健插入式估计量的关系]
再考虑反事实平均结局 $\E[Y^{a}]\equiv\psi_a$，并假定可识别条件成立。于是 $\psi_a=\E[b(a,L)]$，其中 $b(a,L)=\E[Y\given A=a,L]$；同时 $\psi_a=\E\bigl[I(A=a)Y/f(a\given L)\bigr]$。如技术注记 13.2 所述，反事实平均结局 $\E[Y^{a}]\equiv\psi_a$ 的增广 IP 加权（AIPW）估计量 $\hat\psi_{a,AIPW}$ 为
\begin{align*}
&P_n\left[\frac{I(A=a)Y}{\hat f(a\given L)}-\left(\frac{I(A=a)}{\hat f(a\given L)}-1\right)\hat b(a,L)\right]\\
&\qquad=P_n\left[\hat b(a,L)+I(A=a)\left\{Y-\hat b(A,L)\right\}/\hat f(A\given L)\right],
\end{align*}
其中 $P_n[H]\equiv\frac{1}{n}\sum_{i=1}^{n}H_i$ 对任意 $H$ 成立，$\hat f(a\given L)$ 与 $\hat b(a,L)$ 分别是 $f(a\given L)$ 与 $b(a,L)$ 的估计量（Robins et al. 1994, Robins and Ritov 1997）。该估计量是双重稳健的，因为 (i) 若 $\hat f(a\given L)$ 一致，则上述等式左端依概率收敛到 $\psi_a=\E\bigl[I(A=a)Y/f(a\given L)\bigr]$；(ii) 若 $\hat b(a,L)$ 一致，则等式右端收敛到 $\psi_a=\E[b(a,L)]$。由此可知，平均因果效应 $\E[Y^{1}]-\E[Y^{0}]=\psi_1-\psi_0$ 的 AIPW 估计量 $\hat\psi_{1,AIPW}-\hat\psi_{0,AIPW}$ 是双重稳健的：只要 (i) $\hat f(a=1\given L)=1-\hat f(a=0\given L)$ 一致，或 (ii) $\hat b(1,L)$ 与 $\hat b(0,L)$ 都一致，它就是一致的。

既然有了双重稳健的 AIPW 估计量，如何得到精要 13.2 的双重稳健插入式估计量？由上述等式的右端，我们有
\begin{align*}
\hat\psi_{1,AIPW}-\hat\psi_{0,AIPW}
&=P_n[\hat b(1,L)]-P_n[\hat b(0,L)]\\
&\quad-P_n\left[\frac{Y-\hat b(A,L)}{\hat f(A\given L)}\left\{I(A=1)-I(A=0)\right\}\right].
\end{align*}
若我们要的是双重稳健插入式估计量 $P_n[\hat b(1,L)]-P_n[\hat b(0,L)]$，就需要在每一个样本中都有
\[
P_n\left[\frac{Y-\hat b(A,L)}{\hat f(A\given L)}\left\{I(A=1)-I(A=0)\right\}\right]=0.
\]
上述方程成立的条件是：$\hat b(A,L)=b(A,L;\hat\beta,\hat\theta)$ 是模型
\[
\E[Y\given A,L]=b(A,L;\beta,\theta)=\phi\left\{m(A,L;\beta)+\theta\frac{\left\{I(A=1)-I(A=0)\right\}}{\hat f(A\given L)}\right\}
\]
的迭代重加权最小二乘（IRLS）估计，其中 $\phi$ 是标准联系函数（如对数、logit 或线性联系）的逆函数。这是因为上述方程正是对应于参数 $\theta$ 的得分方程（Robins 1999, Bang and Robins 2005, Scharfstein et al. 1999）。所得的插入式估计量正是精要 13.2 的估计量。该估计量是一种基于目标最小损失的估计量（targeted minimum loss-based estimator，TMLE），也称目标最大似然估计量；而
\[
\frac{\{I(A=1)-I(A=0)\}}{\hat f(A\given L)}=\frac{A}{\hat\pi(L)}-\frac{(1-A)}{(1-\hat\pi(L))}
\]
就是 van der Laan 与 Rubin (2006) 后来命名的“聪明协变量”（clever covariate）。

能保证上述方程成立的模型选择不止一种。例如，可以把上述模型中的 $\theta$ 项换成和式 $\theta_1\frac{A}{\hat\pi(L)}+\theta_2\frac{(1-A)}{(1-\hat\pi(L))}$，并同时估计 $\theta_1$ 与 $\theta_2$（Scharfstein et al. 1999, Bang and Robins 2005）。后一个估计量也是 TMLE，但如今它同时也是 $\E[Y^{a=1}]$ 与 $\E[Y^{a=0}]$ 的双重稳健插入式估计量，而 $P_n[\hat b(1,L)]-P_n[\hat b(0,L)]$ 仍是平均治疗效应的双重稳健估计量。双聪明协变量模型相对于单聪明协变量模型的一个优点是：$P_n[\hat b(1,L)]$ 与 $P_n[\hat b(0,L)]$ 现在也分别是 $\E[Y^{a=1}]$ 与 $\E[Y^{a=0}]$ 的双重稳健插入式估计量。
\end{techpoint}
\renewcommand{\TPname}{技术注记}
